\chapimg{fig/at/op_at19.jpg}
\chapterlead{\subhd{After this chapter you can} read every figure of this part of the results atlas, name the quantity on each axis, write the equation from which the figure is computed, say which experiment supplies the data, and judge what the figure does and does not show.}
\chapter{The Atlas of Results I: the System, the Soil Pores and the Coupled Fluxes}\label{chap:19}
\markboth{The Atlas of Results I: the System, the Soil Pores and the Coupled Fluxes}{The Atlas of Results I: the System, the Soil Pores and the Coupled Fluxes}
Every figure of this chapter is shown with the same five items, in the same order: \emph{what it shows}, \emph{the equation behind it} with the meaning of its symbols, \emph{how it is derived} from the Daisy outputs, \emph{the data} it uses, and \emph{how to read it}. The experiments are named E0 to E16 as in the earlier chapters; ``Daisy output used'' names the columns of the run table that the figure draws on. A figure is read against its equation: the axes are the symbols of the equation, and a feature of the curves that none of them explains is a finding to be checked, not a result to be reported.
\banner{The system, the forcing and the baseline}
\section{Feedback loops with data-derived effect sizes}
\fig{at/a19_01.jpg}{1.0}{Causal-loop diagram; arrow labels are regression slopes from E11, E14, E12, E1/E10 and E13.}\label{fig:a19_1}
\par\noindent\textbf{\color{khmag}\emph{What it shows.}} Causal-loop diagram; arrow labels are regression slopes from E11, E14, E12, E1/E10 and E13.
\par\noindent\textbf{\color{khmag}\emph{The equation.}} The figure is computed from
\[b=\dfrac{\sum_i (x_i-\bar x)(y_i-\bar y)}{\sum_i (x_i-\bar x)^2}\]
\par\noindent where $x_i,y_i$ the factor level and the output of run $i$; $\bar x,\bar y$ their means; $b$ the least-squares slope.
\par\noindent\textbf{\color{khmag}\emph{How it is derived.}} Ten arrows are ten least-squares slopes, each from the runs of one experiment with the other factors fixed (effect\_sizes.csv).
\par\noindent\textbf{\color{khmag}\emph{The data.}} experiments and sources: E1,E10-E14; Daisy output used: irrigation\_mm, ET\_mm, soil\_evap\_mm, drainage\_1m\_mm, N\_leached\_kg\_ha, yield\_Mg\_ha.
\section{Experiment design matrix}
\fig{at/a19_02.jpg}{1.0}{Number of levels of each factor per experiment (from the run names) and the number of simulations; red = runs excluded for solver failure.}\label{fig:a19_2}
\par\noindent\textbf{\color{khmag}\emph{What it shows.}} Number of levels of each factor per experiment (from the run names) and the number of simulations; red = runs excluded for solver failure.
\par\noindent\textbf{\color{khmag}\emph{The equation.}} The figure is computed from
\[N=\sum_{e=0}^{16} n_e=10\,190,\qquad N_{\mathrm{valid}}=10\,041\]
\par\noindent where $n_e$ number of runs of experiment $e$ (E0 to E16); $N$ total runs; 149 solver failures are excluded.
\par\noindent\textbf{\color{khmag}\emph{How it is derived.}} The matrix lists which factor varies in which experiment; counts are taken from the run table.
\par\noindent\textbf{\color{khmag}\emph{The data.}} experiments and sources: All; Daisy output used: run table columns experiment, valid.
\section{Climate forcing}
\fig{at/a19_03.jpg}{1.0}{Annual rain, reference ET and temperature for the two climates used (Taastrup: E0--E5, E7--E16; Bologna: E6).}\label{fig:a19_3}
\par\noindent\textbf{\color{khmag}\emph{What it shows.}} Annual rain, reference ET and temperature for the two climates used (Taastrup: E0--E5, E7--E16; Bologna: E6).
\par\noindent\textbf{\color{khmag}\emph{The equation.}} The figure is computed from
\[A=\overline{\left(ET_0/P\right)}\]
\par\noindent where $ET_0$ annual reference evapotranspiration (mm); $P$ annual rain (mm); $A$ mean aridity index.
\par\noindent\textbf{\color{khmag}\emph{How it is derived.}} Annual rain, reference ET and temperature are read from the two weather files (Taastrup and Bologna).
\par\noindent\textbf{\color{khmag}\emph{The data.}} experiments and sources: Weather files, E0; Daisy output used: weather files (rain, reference ET, temperature).
\section{Rainfed yields per year and soil}
\fig{at/a19_04.jpg}{1.0}{Harvest-by-harvest yields from the E0 baselines (re-run for per-year output; matches the user's runs within 0.4 \%).}\label{fig:a19_4}
\par\noindent\textbf{\color{khmag}\emph{What it shows.}} Harvest-by-harvest yields from the E0 baselines (re-run for per-year output; matches the user's runs within 0.4 \%).
\par\noindent\textbf{\color{khmag}\emph{The equation.}} The figure is computed from
\[Y_{s,y}=\mathrm{sorg\_DM}_{s,y}\]
\par\noindent where $Y_{s,y}$ dry-matter yield of soil $s$ in harvest year $y$.
\par\noindent\textbf{\color{khmag}\emph{How it is derived.}} Harvest-by-harvest yields of the ten rainfed baselines, re-run for per-year output.
\par\noindent\textbf{\color{khmag}\emph{The data.}} experiments and sources: E0; Daisy output used: Annual-Harvest.dlf: sorg\_DM.
\section{Rainfed water balance per soil}
\fig{at/a19_05.jpg}{1.0}{Mean annual partitioning of rain in the 10 rainfed baselines.}\label{fig:a19_5}
\par\noindent\textbf{\color{khmag}\emph{What it shows.}} Mean annual partitioning of rain in the 10 rainfed baselines.
\par\noindent\textbf{\color{khmag}\emph{The equation.}} The figure is computed from
\[P+I=T+E_s+E_c+E_o+D+\Delta S+R,\qquad E_o=ET-T-E_s-E_c\]
\par\noindent where $P$ rain; $I$ irrigation; $T$ transpiration; $E_s$ soil evaporation; $E_c$ canopy evaporation; $E_o$ evaporation of ponded water, snow and litter or mulch; $D$ drainage across 1 m; $\Delta S$ storage change; $R$ runoff; $ET$ actual evapotranspiration (all mm yr$^{-1}$).
\par\noindent\textbf{\color{khmag}\emph{How it is derived.}} Mean annual partition of rain in the ten rainfed baselines; $E_o$ is the closing term.
\par\noindent\textbf{\color{khmag}\emph{The data.}} experiments and sources: E0; Daisy output used: Annual-FWater\_100cm.dlf, Annual-SWater.dlf.
\banner{The pore scale and the structure of the soils}
\section{Retention curves of the 10 soils}
\fig{at/a19_06.jpg}{1.0}{van Genuchten $\theta$(h) with field capacity (pF 2), the -300 cm trigger and wilting point (pF 4.2).}\label{fig:a19_6}
\par\noindent\textbf{\color{khmag}\emph{What it shows.}} van Genuchten $\theta$(h) with field capacity (pF 2), the -300 cm trigger and wilting point (pF 4.2).
\par\noindent\textbf{\color{khmag}\emph{The equation.}} The figure is computed from
\[\theta(h)=\theta_r+\dfrac{\theta_s-\theta_r}{\left[1+(\alpha|h|)^{n}\right]^{m}},\quad m=1-\dfrac1n\]
\par\noindent where $\theta$ volumetric water content (cm$^3$ cm$^{-3}$); $h$ suction head (cm); $\theta_r,\theta_s$ residual and saturated water contents; $\alpha$ (cm$^{-1}$) and $n$ van Genuchten shape parameters.
\par\noindent\textbf{\color{khmag}\emph{How it is derived.}} The van Genuchten curve of each soil and horizon, with field capacity at pF 2, the trigger at 300 cm and the wilting point at pF 4.2 ($\mathrm{pF}=\log_{10}|h|$).
\par\noindent\textbf{\color{khmag}\emph{The data.}} experiments and sources: Soil files; Daisy output used: soil files (parameters only).
\section{Pore-size distribution}
\fig{at/a19_07.jpg}{1.0}{Derivative of the retention curve (-d$\theta$/dpF) for each soil; top axis converts suction to equivalent pore radius.}\label{fig:a19_7}
\par\noindent\textbf{\color{khmag}\emph{What it shows.}} Derivative of the retention curve (-d$\theta$/dpF) for each soil; top axis converts suction to equivalent pore radius.
\par\noindent\textbf{\color{khmag}\emph{The equation.}} The figure is computed from
\[\text{density}=-\dfrac{d\theta}{d\,\mathrm{pF}}\]
\[r\,(\mu\mathrm{m})=\dfrac{0.149}{h}\times10^{4}\]
\par\noindent where pF$=\log_{10}|h|$; $h$ suction (cm); $r$ equivalent pore radius from the Young--Laplace law $r=2\sigma\cos\gamma/(\rho g h)=0.149/h$ cm.
\par\noindent\textbf{\color{khmag}\emph{How it is derived.}} The derivative of the retention curve is the pore-volume density; the top axis converts suction to pore radius.
\par\noindent\textbf{\color{khmag}\emph{The data.}} experiments and sources: Soil files; Daisy output used: soil files (parameters only).
\section{Unsaturated conductivity}
\fig{at/a19_08.jpg}{1.0}{Mualem--van Genuchten K(Se) and K(h) (Ap horizon), mm/day; dotted line = trigger.}\label{fig:a19_8}
\par\noindent\textbf{\color{khmag}\emph{What it shows.}} Mualem--van Genuchten K(Se) and K(h) (Ap horizon), mm/day; dotted line = trigger.
\par\noindent\textbf{\color{khmag}\emph{The equation.}} The figure is computed from
\[K(S_e)=K_s\,S_e^{\,\ell}\left[1-\left(1-S_e^{1/m}\right)^{m}\right]^{2},\quad S_e=\left[1+(\alpha|h|)^{n}\right]^{-m}\]
\par\noindent where $K$ hydraulic conductivity; $K_s$ saturated conductivity; $S_e$ effective saturation; $\ell$ pore-connectivity exponent; $m=1-1/n$.
\par\noindent\textbf{\color{khmag}\emph{How it is derived.}} Mualem--van Genuchten conductivity of the Ap horizon, converted to mm day$^{-1}$ (cm h$^{-1}\times240$).
\par\noindent\textbf{\color{khmag}\emph{The data.}} experiments and sources: Soil files; Daisy output used: soil files (parameters only).
\section{Available water, air entry and Ksat}
\fig{at/a19_09.jpg}{1.0}{Soil-physical indicators computed from the hydraulic parameters of each soil.}\label{fig:a19_9}
\par\noindent\textbf{\color{khmag}\emph{What it shows.}} Soil-physical indicators computed from the hydraulic parameters of each soil.
\par\noindent\textbf{\color{khmag}\emph{The equation.}} The figure is computed from
\[\mathrm{TAW}=\sum_{k\in\{Ap,Bt\}}\bigl[\theta_k(100)-\theta_k(15849)\bigr]\,\Delta z_k\times10\]
\[h_{\mathrm{ae}}=1/\alpha\]
\par\noindent where TAW total available water 0--100 cm (mm); $\theta_k(h)$ van Genuchten water content of horizon $k$; $\Delta z_k$ thickness (cm; Ap 25 cm, Bt 75 cm); $h_{\mathrm{ae}}$ air-entry suction.
\par\noindent\textbf{\color{khmag}\emph{How it is derived.}} Plant-available water from field capacity (100 cm) to wilting point (15\,849 cm), air entry from $1/\alpha$, and $K_s$ from the soil file.
\par\noindent\textbf{\color{khmag}\emph{The data.}} experiments and sources: Soil files; Daisy output used: soil files (parameters only).
\section{Topsoil conductivity and water content over time}
\fig{at/a19_10.jpg}{1.0}{Daily log10 K and $\theta$ at 3.75 cm (30-day rolling statistics), E5 re-runs.}\label{fig:a19_10}
\par\noindent\textbf{\color{khmag}\emph{What it shows.}} Daily log10 K and $\theta$ at 3.75 cm (30-day rolling statistics), E5 re-runs.
\par\noindent\textbf{\color{khmag}\emph{The equation.}} The figure is computed from
\[K(t)=K\bigl(h(t)\bigr),\quad \bar x_{30}(t)=\dfrac1{30}\sum_{j=0}^{29}x(t-j)\]
\par\noindent where $K(h)$ Mualem--van Genuchten conductivity at the daily suction $h$; $\bar x_{30}$ 30-day rolling mean of $\log_{10}K$ or $\theta$.
\par\noindent\textbf{\color{khmag}\emph{How it is derived.}} Daily log K and water content at 3.75 cm, smoothed with a 30-day rolling mean.
\par\noindent\textbf{\color{khmag}\emph{The data.}} experiments and sources: E5 daily; Daisy output used: profile logs: h @ -3.75, Theta @ -3.75.
\section{Structural degradation rate effects}
\fig{at/a19_11.jpg}{1.0}{E5: change in yield, drainage and irrigation relative to tc = 0.001/day, per tillage system.}\label{fig:a19_11}
\par\noindent\textbf{\color{khmag}\emph{What it shows.}} E5: change in yield, drainage and irrigation relative to tc = 0.001/day, per tillage system.
\par\noindent\textbf{\color{khmag}\emph{The equation.}} The figure is computed from
\[\Delta_k=\overline{x_k-x_{k,b}}\]
\par\noindent where $x_k$ yield, drainage or irrigation at consolidation rate $k$; $x_{k,b}$ the same at the baseline rate $t_c=0.001$ d$^{-1}$; the bar the mean over soils.
\par\noindent\textbf{\color{khmag}\emph{How it is derived.}} The change in each output relative to the baseline structural-degradation rate, per tillage system.
\par\noindent\textbf{\color{khmag}\emph{The data.}} experiments and sources: E5; Daisy output used: yield\_Mg\_ha, drainage\_1m\_mm, irrigation\_mm (E5).
\section{Tillage $\times$ water-regime interaction}
\fig{at/a19_12.jpg}{1.0}{E5 means over the five consolidation rates.}\label{fig:a19_12}
\par\noindent\textbf{\color{khmag}\emph{What it shows.}} E5 means over the five consolidation rates.
\par\noindent\textbf{\color{khmag}\emph{The equation.}} The figure is computed from
\[\bar y_{s,t,w}=\dfrac1{n}\sum_{r\in(s,t,w)} y_r\]
\par\noindent where $y_r$ output of run $r$; $s$ soil; $t$ tillage system; $w$ water regime (irrigated or rainfed); mean over the five consolidation rates.
\par\noindent\textbf{\color{khmag}\emph{How it is derived.}} Means of E5 outputs by soil, tillage and water regime.
\par\noindent\textbf{\color{khmag}\emph{The data.}} experiments and sources: E5; Daisy output used: yield\_Mg\_ha, irrigation\_mm, drainage\_1m\_mm (E5).
\section{Wetting--drying loops}
\fig{at/a19_13.jpg}{1.0}{Daily $\theta$ vs pF at 3.75 cm through the 1994 season for each tillage system.}\label{fig:a19_13}
\par\noindent\textbf{\color{khmag}\emph{What it shows.}} Daily $\theta$ vs pF at 3.75 cm through the 1994 season for each tillage system.
\par\noindent\textbf{\color{khmag}\emph{The equation.}} The figure is computed from
\[\mathrm{pF}(t)=\log_{10}\bigl(-h(t)\bigr)\ \ \text{plotted against}\ \ \theta(t)\]
\par\noindent where $h(t)$ daily suction and $\theta(t)$ daily water content at 3.75 cm during 1994.
\par\noindent\textbf{\color{khmag}\emph{How it is derived.}} Hysteresis-like loops appear because the same pF has a different $\theta$ on the wetting and drying branch.
\par\noindent\textbf{\color{khmag}\emph{The data.}} experiments and sources: E5 daily; Daisy output used: profile logs: h @ -3.75, Theta @ -3.75.
\banner{The coupled water fluxes}
\section{Water-flux Sankey, loam vs sand}
\fig{at/a19_14.jpg}{1.0}{Mean annual inputs (rain, irrigation) and outputs (T, E, interception, other evaporation, drainage) from E2 at the course trigger.}\label{fig:a19_14}
\par\noindent\textbf{\color{khmag}\emph{What it shows.}} Mean annual inputs (rain, irrigation) and outputs (T, E, interception, other evaporation, drainage) from E2 at the course trigger.
\par\noindent\textbf{\color{khmag}\emph{The equation.}} The figure is computed from
\[P+I=T+E_s+E_c+E_o+D+\Delta S+R,\qquad E_o=ET-T-E_s-E_c\]
\[R_{\mathrm{res}}=P+I-(T+E_s+E_c+E_o+D)\]
\par\noindent where $P$ rain; $I$ irrigation; $T$ transpiration; $E_s$ soil evaporation; $E_c$ canopy evaporation; $E_o$ evaporation of ponded water, snow and litter or mulch; $D$ drainage across 1 m; $\Delta S$ storage change; $R$ runoff; $ET$ actual evapotranspiration (all mm yr$^{-1}$); $R_{\mathrm{res}}$ storage or runoff, drawn when larger than 3 mm.
\par\noindent\textbf{\color{khmag}\emph{How it is derived.}} Inputs (rain, irrigation) and outputs of the 20-year mean annual balance of the runs E2 at the $-300$ cm trigger on coarse loam and coarse sand.
\par\noindent\textbf{\color{khmag}\emph{The data.}} experiments and sources: E2; Daisy output used: rain\_mm, irrigation\_mm, transpiration\_mm, soil\_evap\_mm, canopy\_evap\_mm, drainage\_1m\_mm.
\section{Seasonal E/T partitioning}
\fig{at/a19_15.jpg}{1.0}{Daily transpiration and soil evaporation (7-day means) with LAI, rainfed vs irrigated vs mulched.}\label{fig:a19_15}
\par\noindent\textbf{\color{khmag}\emph{What it shows.}} Daily transpiration and soil evaporation (7-day means) with LAI, rainfed vs irrigated vs mulched.
\par\noindent\textbf{\color{khmag}\emph{The equation.}} The figure is computed from
\[\text{share}_{E}=\dfrac{\sum E_s}{\sum E_s+\sum T}\]
\par\noindent where $E_s$ daily soil evaporation; $T$ daily transpiration (both 7-day means); sums over the season.
\par\noindent\textbf{\color{khmag}\emph{How it is derived.}} The soil-evaporation share of evaporation plus transpiration appears in each panel title.
\par\noindent\textbf{\color{khmag}\emph{The data.}} experiments and sources: E0,E2,E7 daily; Daisy output used: Daily-SWater.dlf: Actual transpiration, Evaporation of soil water; LAI.
\section{Daily drainage vs rain events}
\fig{at/a19_16.jpg}{1.0}{Daily rain and irrigation (bars, top) and drainage at 1 m (shaded, inverted) in the driest and wettest years.}\label{fig:a19_16}
\par\noindent\textbf{\color{khmag}\emph{What it shows.}} Daily rain and irrigation (bars, top) and drainage at 1 m (shaded, inverted) in the driest and wettest years.
\par\noindent\textbf{\color{khmag}\emph{The equation.}} The figure is computed from
\[D_d=-q(z{=}100\ \mathrm{cm},d)\]
\par\noindent where $D_d$ daily drainage across 1 m on day $d$; $q$ downward water flux.
\par\noindent\textbf{\color{khmag}\emph{How it is derived.}} Daily rain and irrigation (bars) and drainage at 1 m (shaded) for the driest and wettest years.
\par\noindent\textbf{\color{khmag}\emph{The data.}} experiments and sources: E2 daily; Daisy output used: Daily-SWater.dlf: rain, Surface irrigation; Daily-FWater\_100cm.dlf.
\section{Cumulative fluxes by method}
\fig{at/a19_17.jpg}{1.0}{Cumulative irrigation, drainage and canopy evaporation 1980--2000 for four methods.}\label{fig:a19_17}
\par\noindent\textbf{\color{khmag}\emph{What it shows.}} Cumulative irrigation, drainage and canopy evaporation 1980--2000 for four methods.
\par\noindent\textbf{\color{khmag}\emph{The equation.}} The figure is computed from
\[C_x(t)=\sum_{\tau\le t} x_\tau\]
\par\noindent where $x_\tau$ daily irrigation, drainage or canopy evaporation; $C_x(t)$ cumulative total from 1980 to day $t$.
\par\noindent\textbf{\color{khmag}\emph{How it is derived.}} Cumulative curves of four irrigation methods on coarse loam.
\par\noindent\textbf{\color{khmag}\emph{The data.}} experiments and sources: E4 daily; Daisy output used: Daily irrigation, drainage, Evaporation from canopy.
\section{Partitioning by method $\times$ soil}
\fig{at/a19_18.jpg}{1.0}{Stacked mean annual fluxes for 10 methods in all 10 soils; dots = irrigation.}\label{fig:a19_18}
\par\noindent\textbf{\color{khmag}\emph{What it shows.}} Stacked mean annual fluxes for 10 methods in all 10 soils; dots = irrigation.
\par\noindent\textbf{\color{khmag}\emph{The equation.}} The figure is computed from
\[P+I=T+E_s+E_c+E_o+D+\Delta S+R,\qquad E_o=ET-T-E_s-E_c\]
\par\noindent where $P$ rain; $I$ irrigation; $T$ transpiration; $E_s$ soil evaporation; $E_c$ canopy evaporation; $E_o$ evaporation of ponded water, snow and litter or mulch; $D$ drainage across 1 m; $\Delta S$ storage change; $R$ runoff; $ET$ actual evapotranspiration (all mm yr$^{-1}$).
\par\noindent\textbf{\color{khmag}\emph{How it is derived.}} Mean annual fluxes stacked for ten methods in ten soils; dots mark irrigation.
\par\noindent\textbf{\color{khmag}\emph{The data.}} experiments and sources: E4; Daisy output used: E4 annual outputs.
\section{Drip depth vs drainage and evaporation}
\fig{at/a19_19.jpg}{1.0}{E4 lines per soil; E15 averaged over the 10 amounts per rooting depth.}\label{fig:a19_19}
\par\noindent\textbf{\color{khmag}\emph{What it shows.}} E4 lines per soil; E15 averaged over the 10 amounts per rooting depth.
\par\noindent\textbf{\color{khmag}\emph{The equation.}} The figure is computed from
\[\mathrm{DPF}=\dfrac{D}{P+I}\]
\[z_{\mathrm{mid}}=\dfrac{z_1+z_2}2\]
\par\noindent where $D$ drainage across 1 m depth (mm yr$^{-1}$); $P$ rain; $I$ irrigation; DPF deep-percolation fraction; $z_1,z_2$ top and bottom depth of the drip band (cm); $z_{\mathrm{mid}}$ its mid-depth.
\par\noindent\textbf{\color{khmag}\emph{How it is derived.}} DPF and soil evaporation plotted against the mid-depth of the drip band; E15 averaged over the ten amounts per rooting depth.
\par\noindent\textbf{\color{khmag}\emph{The data.}} experiments and sources: E4,E15; Daisy output used: drainage\_1m\_mm, rain\_mm, irrigation\_mm, soil\_evap\_mm (E4, E15).
\section{Method $\times$ rooting depth $\times$ amount}
\fig{at/a19_20.jpg}{1.0}{IE for 1,000 runs; each panel one application method.}\label{fig:a19_20}
\par\noindent\textbf{\color{khmag}\emph{What it shows.}} IE for 1,000 runs; each panel one application method.
\par\noindent\textbf{\color{khmag}\emph{The equation.}} The figure is computed from
\[\mathrm{IE}=\dfrac{1000\,Y}{10\,I}\]
\par\noindent where $Y$ yield (Mg ha$^{-1}$); $I$ irrigation (mm yr$^{-1}$); IE irrigation efficiency (kg m$^{-3}$); kept only where $I\ge 20$ mm yr$^{-1}$.
\par\noindent\textbf{\color{khmag}\emph{How it is derived.}} IE of 1,000 runs, one panel per application method; colour scale from the 2nd to the 98th percentile.
\par\noindent\textbf{\color{khmag}\emph{The data.}} experiments and sources: E15; Daisy output used: yield\_Mg\_ha, irrigation\_mm (E15).
\section{Annual drainage vs water input}
\fig{at/a19_21.jpg}{1.0}{One point per year; lines = linear fits per soil.}\label{fig:a19_21}
\par\noindent\textbf{\color{khmag}\emph{What it shows.}} One point per year; lines = linear fits per soil.
\par\noindent\textbf{\color{khmag}\emph{The equation.}} The figure is computed from
\[D=a+b\,P,\qquad D=a'+b'\,(P+I)\]
\par\noindent where $D$ annual drainage; $P$ annual rain; $I$ annual irrigation; $a,b,a',b'$ least-squares intercepts and slopes per soil.
\par\noindent\textbf{\color{khmag}\emph{How it is derived.}} One point per year; the lines are linear fits per soil.
\par\noindent\textbf{\color{khmag}\emph{The data.}} experiments and sources: E0,E2 per year; Daisy output used: Annual-FWater\_100cm.dlf: Precipitation, Irrigation, Matrix percolation.
\section{Profile pressure-head evolution}
\fig{at/a19_22.jpg}{1.0}{Depth $\times$ date pF with the trigger iso-line; dashed line = sensor depth.}\label{fig:a19_22}
\par\noindent\textbf{\color{khmag}\emph{What it shows.}} Depth $\times$ date pF with the trigger iso-line; dashed line = sensor depth.
\par\noindent\textbf{\color{khmag}\emph{The equation.}} The figure is computed from
\[\mathrm{pF}(z,t)=\log_{10}\bigl(-h(z,t)\bigr),\qquad \text{contour at } \mathrm{pF}=\log_{10}300\]
\par\noindent where $h(z,t)$ suction at depth $z$ and date $t$; the black line is the $-300$ cm trigger iso-line.
\par\noindent\textbf{\color{khmag}\emph{How it is derived.}} Depth by date map of pF in 1994 with the trigger iso-line and the sensor depth.
\par\noindent\textbf{\color{khmag}\emph{The data.}} experiments and sources: E2 daily; Daisy output used: profile logs: h @ depth.
\section{Year-to-year variability}
\fig{at/a19_23.jpg}{1.0}{Boxplots across the 20 harvest years for each trigger and method.}\label{fig:a19_23}
\par\noindent\textbf{\color{khmag}\emph{What it shows.}} Boxplots across the 20 harvest years for each trigger and method.
\par\noindent\textbf{\color{khmag}\emph{The equation.}} The figure is computed from
\[\text{box}=(Q_{25},Q_{50},Q_{75}),\ \text{whiskers}=1.5\,\mathrm{IQR}\]
\par\noindent where $Q_p$ percentile over the 20 harvest years; IQR interquartile range.
\par\noindent\textbf{\color{khmag}\emph{How it is derived.}} Boxplots of the yearly values for each trigger and each method.
\par\noindent\textbf{\color{khmag}\emph{The data.}} experiments and sources: E2,E4 per year; Daisy output used: per-year Annual-Harvest.dlf, Annual-FWater\_100cm.dlf.
\begin{itemize}
\item A figure is a statement about one equation and one data set; both are named under it.
\item The ten soils are ordered by the shape parameter $n$ of the van Genuchten curve; many figures are read along that order.
\item IE and IWUE are shown only where the irrigation is at least 20 mm per year.
\end{itemize}
\backhead{EXERCISES}{khred}\begin{exlist}
\item Write the equation from which the figure ``Feedback loops with data-derived effect sizes'' is computed and state what each symbol means. \exans{see the equation under the figure}
\item Write the equation from which the figure ``Experiment design matrix'' is computed and state what each symbol means. \exans{see the equation under the figure}
\item Write the equation from which the figure ``Climate forcing'' is computed and state what each symbol means. \exans{see the equation under the figure}
\item Write the equation from which the figure ``Rainfed yields per year and soil'' is computed and state what each symbol means. \exans{see the equation under the figure}
\item Write the equation from which the figure ``Rainfed water balance per soil'' is computed and state what each symbol means. \exans{see the equation under the figure}
\item Write the equation from which the figure ``Retention curves of the 10 soils'' is computed and state what each symbol means. \exans{see the equation under the figure}
\end{exlist}
\dykband{DO YOU KNOW ?}\begin{exlist}
\item What does the figure ``Feedback loops with data-derived effect sizes'' show, and from which experiment does its data come?
\item What does the figure ``Experiment design matrix'' show, and from which experiment does its data come?
\item What does the figure ``Climate forcing'' show, and from which experiment does its data come?
\item What does the figure ``Rainfed yields per year and soil'' show, and from which experiment does its data come?
\item What does the figure ``Rainfed water balance per soil'' show, and from which experiment does its data come?
\item What does the figure ``Retention curves of the 10 soils'' show, and from which experiment does its data come?
\item What does the figure ``Pore-size distribution'' show, and from which experiment does its data come?
\end{exlist}
\backhead{OBJECTIVE TYPE QUESTIONS}{khred}\begin{exlist}
\item The data of the figure ``Retention curves of the 10 soils'' come from
\opt{E1-E4,E10}{Soil files}{E5 daily}{E5, E0 daily}
\item The data of the figure ``Pore-size distribution'' come from
\opt{Soil files}{E0,E2 per year}{E7,E2}{Pareto set}
\item The data of the figure ``Profile pressure-head evolution'' come from
\opt{E2 daily}{E2 + soil hydraulics}{E1-E7}{E2,E4 per year}
\item The data of the figure ``Annual drainage vs water input'' come from
\opt{E7,E14}{E0,E2 per year}{E4 daily}{E0,E2 daily}
\item The data of the figure ``Year-to-year variability'' come from
\opt{E2,E4 per year}{E1-E7}{E5, E0 daily}{E0,E2 per year}
\end{exlist}
\objkey{1.\ (b),\quad 2.\ (a),\quad 3.\ (a),\quad 4.\ (b),\quad 5.\ (a)}
\swaflvideo{19}
